% !PRINTED 463
\weberpage{463}
%<<CONT
\[
\Phi(t) = (t - \Theta)(t - \Theta_1) \ldots (t - \Theta_{n-1}),
\]
so ist
\[
\Phi(t) \left( \frac{\omega}{t - \Theta} + \frac{\omega_1}{t - \Theta_1} + \cdots + \frac{\omega_{n-1}}{t - \Theta_{n-1}} \right) = \Psi(t)
\]
eine ganze rationale Function $n - 1^{\mathrm{ten}}$ Grades von $t$,
deren Co\"efficienten Zahlen in $\Omega$ sind, und daraus ergiebt
sich, wenn man $t = \Theta$ setzt,
\[
\omega = \frac{\Psi(\Theta)}{\Phi'(\Theta)},
\]
worin $\Phi'(\Theta)$ von Null verschieden ist.

\gesperrt{Es ist hiernach der Körper $\Omega(\Theta)$ mit dem Körper
$\Omega(\alpha)$ identisch.}

\gesperrt{4. Ist $\Theta$ nicht primitiv, so kann nicht jede Grösse
in $\Omega(\alpha)$ rational durch $\Theta$ ausgedrückt werden. Der
Körper $\Omega(\Theta)$ ist ein Theiler des Körpers $\Omega(\alpha)$
und der Grad von $\Omega(\Theta)$ ist ein Theiler des Grades von
$\Omega(\alpha)$.}

Denn jede Zahl des Körpers $\Omega(\Theta)$ genügt einer Gleichung
in $\Omega$ vom Grade $n_1$, wenn $n_1$ ein Theiler von $n$ und
kleiner als $n$ ist. Also kann eine primitive Grösse des Körpers
$\Omega(\alpha)$, die einer irreducibeln Gleichung $n^{\mathrm{ten}}$
Grades genügt, nicht in $\Omega(\Theta)$ enthalten sein.

\gesperrt{Der Körper $\Omega(\alpha)$ heisst primitiv, wenn er ausser den
Grössen in $\Omega$ keine imprimitiven Grössen enthält, imprimitiv,
wenn er noch andere imprimitive Grössen enthält.}

Aus dieser Definition ergiebt sich zunächst, dass ein Körper,
dessen Grad eine Primzahl ist, nothwendig primitiv ist; denn eine
imprimitive Grösse $\Theta$ in $\Omega(\alpha)$ genügt einer
Gleichung, deren Grad ein von $n$ verschiedener Theiler des Grades
$n$ von $\Omega(\alpha)$ ist. Wenn dieser Theiler gleich 1 ist, so
ist $\Theta$ in $\Omega$ enthalten.

Wir wollen jetzt noch einige der wichtigsten Eigenschaften der
imprimitiven Körper kennen lernen.

Hat $\Omega(\alpha)$ einen von $\Omega$ verschiedenen Theiler
$\Omega'$, der seinerseits $\Omega$ als Theiler enthält (ist also
$\Omega'$ Körper über $\Omega$), und ist $\beta$ eine Grösse, die
dem Körper $\Omega'$, aber nicht $\Omega$ angehört, so ist der
Körper $\Omega(\beta)$ ein algebraischer Körper über $\Omega$ und
zugleich ein Theiler von $\Omega'$ und von $\Omega(\alpha)$, und
nach dem, was wir vorhin
%CONT>>
